\documentclass[10pt,a4paper]{amsart}
\usepackage[margin=19mm]{geometry}
\usepackage{amsmath,amssymb,mathtools}
\usepackage[scr=rsfs,cal=euler]{mathalfa}
\usepackage{xfrac}
\usepackage[unicode,hidelinks,bookmarks=false]{hyperref}
\newcommand{\ie}{i.e.\ }
\newcommand{\cf}{cf.\ }
\newcommand{\resp}{respectively\ }
\newcommand{\Subsection}{\subsection}
\providecommand{\nobreakdash}{\nobreak\hbox{-}}
\setlength{\emergencystretch}{2em}
\multlinegap0pt
\usepackage{mathtools}
\usepackage{amssymb}
\usepackage{tikz}
\newtheorem{prop}{Proposition}
\newtheorem{pf}{Pf}
\usepackage{cleveref}
\crefname{prop}{Proposition}{Propositions}
\crefname{pf}{Pf}{Pfs}
\newcommand{\Del}{\Delta}
\newcommand{\Gam}{\Gamma}
\newcommand{\calM}{\mathcal{M}}
\newcommand{\de}{\delta}
\newcommand{\frakp}{\mathfrak{p}}
\newcommand{\gam}{\gamma}
\newcommand{\sig}{\sigma}
\newcommand{\wh}{\widehat}
\title{Finite groups and complex projective surfaces}
\author{Alexander Lubotzky, Matthew Stover}
\date{Source: arXiv:2512.19505v1 (CC BY 4.0)}
\begin{document}
\maketitle

\noindent\emph{Source and licence.} Finite groups and complex projective surfaces. Version arXiv:2512.19505v1 (CC BY 4.0), distributed by arXiv under the Creative Commons Attribution 4.0 International licence, http://creativecommons.org/licenses/by/4.0/, as recorded in the arXiv licence metadata for this exact version and shown on the abstract page https://arxiv.org/abs/2512.19505v1. Full licence notice: \texttt{licenses/arxiv-2512.19505-CC-BY-4.0.txt}.\par\smallskip

\noindent\textit{Editorial scope.} Counted unit: the article's prop prop:FinalNontrivialG (source0.tex lines 816--818). The article's complete original proof of it is delivered with it (source0.tex lines 821--840, one \texttt{proof} environment of 20 lines). No mathematics is written, repaired or re-derived: the statement and the proof are the article's own lines, in the article's own notation, and no retained line is substituted or re-flowed.  Retained uncounted as the source-local context the proof needs: lines 631--633; lines 641--644; lines 646--649; lines 661--663; lines 702--705; lines 728--730; lines 757--764; lines 759--763; lines 771--774.  Nothing was omitted from any retained span, so the package carries no omission record.  No retained line contains a citation, so the wrapper carries no bibliography and no bibliography entry.  No retained line contains a figure environment, an image-inclusion command or drawing code, so no image is delivered and the package declares no asset dependency.  Editorial formatting only, in addition: an AMS \texttt{amsart} preamble that restates the article's own declarations and macros used by the retained lines, namely the environments \texttt{prop}, \texttt{pf} on separate counters, together with the standard packages those lines need.  The retained environments print in this document as Proposition 1, Pf 1 in order of appearance, whereas the article prints the same items with the numbering of its own full document.  The complete native source of the article as distributed in the arXiv e-print for this exact version is bundled unchanged as \texttt{sources/191484/source0.tex}.
\begin{equation}\label{eq:SDef}
S \coloneqq\!\left\{(\psi(t), t)\ :\ t \in T\right\}
\end{equation}

\begin{align}
a_\gam \psi(\rho_{\mathrm{tf}}(\de)) a_\gam^{-1} &= \psi(b_\gam \rho_{\mathrm{tf}}(\de) b_\gam^{-1}) \nonumber \\
&= \psi(\rho_{\mathrm{tf}}(\gam \de \gam^{-1})) \label{eq:1stConj}
\end{align}

\begin{align}
\psi(\rho_{\mathrm{tf}}(\gam \de \gam^{-1})) &= \psi_0( \tau(\gam \de \gam^{-1})) \nonumber \\
&= \psi_0(t_\gam \tau(\de) t_\gam^{-1}) \label{eq:2ndConj}
\end{align}

\begin{equation}\label{eq:CDef}
C \coloneqq\!\left\{\gam \in N_\Gam(\Del)\ :\ \tau(\gam \de \gam^{-1}) = \tau (\de)\ \textrm{for all}\ \de \in \Del \right\}
\end{equation}

\begin{align}
K &\coloneqq \bigcap_{j = 0}^d K_j \label{eq:KDef} \\
k &\coloneqq [\Del : K] \label{eq:kDef}
\end{align}

\begin{prop}\label{prop:NormInD}
If $G$ is nontrivial and $A \in \calM_G$, then $N_\Gam(A) \le D$.
\end{prop}

\begin{prop}\label{prop:FindA}
Suppose $G$ is a nontrivial finite group and $\calM_G$ is defined as in \Cref{prop:NormInD} with respect to $K \triangleleft \Del$ as defined in \Cref{eq:KDef}. After possibly increasing the smallest prime $p_0$ used to define $K$, there exists at least one $A$ in $\calM_G$ so that:
\begin{align}
N_D(A) &= N_\Gam(A) \label{eq:ND=NGam} \\
N_\Del(A) / A &\cong G \label{eq:NDel(A)/A=G} \\
N_D(A) / A &\cong G \times C / M \label{eq:ND(A)/A}
\end{align}
\end{prop}

\begin{align}
N_D(A) &= N_\Gam(A)  \\
N_\Del(A) / A &\cong G  \\
N_D(A) / A &\cong G \times C / M 
\end{align}

\begin{align}
D / M &= \Del / M \times C / M \label{eq:D/M} \\
N_D(A) / M &= N_\Del(A) / M \times C / M \label{eq:ND(A)/M}
\end{align}

\begin{prop}\label{prop:FinalNontrivialG}
Suppose $G$ is a nontrivial finite group, $A \in \calM_G$ satisfies the conclusions of \Cref{prop:FindA}, and $B = AC$ with $C$ as in \Cref{eq:CDef}. Then $N_\Gam(B) = N_\Gam(A)$. In particular, $N_\Gam(B) / B \cong G$.
\end{prop}

\begin{pf}
Since $C$ is normal in $D$, it follows that $N_\Gam(A) \le N_\Gam(B)$. Indeed,
\begin{equation}\label{eq:NDA}
N_D(A) \le N_D(AC) = N_D(B) \le N_\Gam(B)
\end{equation}
by \Cref{prop:NormInD}. The opposite inclusion then proves the proposition, as the second conclusion subsequently follows from \Cref{eq:ND=NGam} and \Cref{eq:ND(A)/A}. Moreover, since $\Del \cap C = M$ by \Cref{eq:D/M} and $M \le A$, it follows that $\Del \cap B = A$. In particular, it suffices to show that $\gam A \gam^{-1} \le \Del$ for all $\gam \in N_\Gam(B)$.

Fix $g \in N_\Gam(B)$, and let $\sig(g) = (g_1, g_2) \in \wh{Q} \times \wh{T}$ denote its reduction modulo $M_\mathrm{tf}(\frakp)$. Recall that the image of $\Del$ in $\wh{Q} \times \wh{T}$ is the group $S$ defined in \Cref{eq:SDef}. Since $C \le N_\Gam(\Del)$ by definition, the image $(a_\gam,b_\gam) \in \wh{Q} \times \wh{T}$ of an element $\gam \in C$ must satisfy \Cref{eq:1stConj} and \Cref{eq:2ndConj} with $t_\gam = 1$. Then $\psi$ factors through the projection $\pi$ onto $T_r$, so $a_\gam = 1$, since $\wh{Q}$ has trivial center. Moreover, the element $b_\gam \in \wh{T}$ must be in the centralizer $Z_{\wh{T}}(T_r)$ of the direct factor $T_r$ of $T$ in order to act trivially through $\pi$.

Now, $ac \in B$ satisfies $\sig(ac) = (\psi(t_a), t_a h_c)$ for some $t_a \in \rho_{\mathrm{tf}}(A)$ and $h_c \in Z_{\wh{T}}(T_r)$. If $a \in A$, then $g a g^{-1} = a^\prime c \in AC$, which gives
\begin{align}
\sig\!\left(g a g^{-1}\right) &= \left(\psi(t^\prime), t^\prime h\right) \nonumber \\
&= \left(g_1 \psi(t) g_1^{-1}, g_2 t g_2^{-1}\right) \label{eq:ACendgame}
\end{align}
for the appropriate $t, t^\prime \in T$ and $h \in Z_{\wh{T}}(T_r)$. Then $g_2 \in \wh{T}$ normalizes $T$, so $t^\prime h \in T$ and then $h \in T$. However, this implies that
\begin{equation}\label{eq:hCentralizes}
h \in \{1\} \times \prod_{\ell \neq r} T_\ell
\end{equation}
in order to centralize the $T_r$ factor, hence $\psi(t^\prime h) = \psi(t^\prime)$ since $\psi$ factors through projection onto $T_r$, meaning precisely that $\sig(g a g^{-1}) \in S$. This proves that $g a g^{-1} \in \Del$ for all $g \in N_\Gam(B)$ and $a \in A$, which is what we needed to prove.
\end{pf}

\end{document}
